\documentclass{exam}
\usepackage{ModeloIFBA}
%\usepackage[alf]{abntex2cite}
\usepackage{multicol}
\usepackage{adjustbox}
\usepackage{multicol}

\setlength\columnsep{30pt}


\renewcommand{\curso}{Tecnologia em Análise e Desenvolvimento de Sistemas }
\renewcommand{\materia}{Programação Orientada a Objetos}
\renewcommand{\turma}{2025.1}
\renewcommand{\professor}{Leandro Costa Souza}
\renewcommand{\avaliacao}{Prova Unidade II - Modelo B}
\renewcommand{\dataavaliacao}{01/09/2025}


\newcommand{\duracao}{1h 40min}
\newcounter{solcount}
\setcounter{solcount}{1}
\newcommand\CC{
    \CorrectChoice\label{sol:\arabic{solcount}}
    \stepcounter{solcount}
}

\newcommand\printmyanswers{
\newpage
\section*{\centering Answers Sheet}
    \foreach \x in {1,...,\thequestion}{
    \noindent
        \x.\ref{sol:\x}\\par
    }
}

\printanswers

\begin{document}
\begin{adjustbox}{max width=\textwidth}
\begin{tabularx}{\textwidth}{rrXrl}
\multirow{5}{*}{\ifbalogo{1.7in}} & \textbf{Curso:} & \multicolumn{3}{l}{ \curso } \\
& \textbf{Matéria:} & \materia \space & \textbf{Turma:} & \small{\turma}  \\
& \textbf{Professor:} & \professor & & \\
& \textbf{Aluno:} & \hrulefill & \textbf{Data:} & \dataavaliacao \\
\end{tabularx}
\end{adjustbox}
\begin{tabularx}{\textwidth}{XcX}
 & \textbf{\avaliacao} &  \ \hline
\end{tabularx}	
	\input{orientacoes}

	\begin{questions}
\begin{multicols}{2}

% --- AMBIENTE E FERRAMENTAS (4 Questões) ---

\question[1] O que a sigla JDK significa no contexto do desenvolvimento Java?
\begin{choices}
 \choice Java Runtime Kit - Kit de Execução Java
 \CC Java Development Kit - Kit de Desenvolvimento Java
 \choice Java Deployment Kit - Kit de Implantação Java
 \choice Java Dynamic Kit - Kit Dinâmico Java
 \choice Java Developer Kernel - Núcleo de Desenvolvimento Java
\end{choices}

\question[1] Qual é a função da JVM (Java Virtual Machine)?
\begin{choices}
 \choice Compilar o código-fonte Java (.java) em bytecode (.class).
 \CC Fornecer um ambiente para executar o bytecode Java em qualquer plataforma.
 \choice Depurar o código Java durante o desenvolvimento.
 \choice Gerenciar as dependências do projeto.
 \choice Escrever código Java automaticamente.
\end{choices}

\question[1] O que o JRE (Java Runtime Environment) inclui?
\begin{choices}
 \choice Apenas o compilador Java (javac).
 \CC A JVM e as bibliotecas padrão do Java, mas não as ferramentas de desenvolvimento.
 \choice O JDK completo.
 \choice Apenas a JVM.
 \choice Ferramentas como o VSCode.
\end{choices}

\question[1] No contexto da JVM, JRE e JDK, qual deles é necessário apenas para executar uma aplicação Java, sem a necessidade de compilá-la?
\begin{choices}
 \choice JVM
 \CC JRE
 \choice JDK
 \choice Javac
 \choice Maven
\end{choices}

% --- FUNDAMENTOS DE JAVA (10 Questões) ---

\question[1] Qual a principal diferença entre tipos primitivos e tipos por referência em Java?
\begin{choices}
 \choice Tipos primitivos são mais rápidos.
 \choice Tipos primitivos armazenam o valor, enquanto tipos por referência armazenam um endereço de memória.
 \choice Tipos primitivos não podem ter métodos.
 \choice Tipos por referência são sempre mutáveis.
 \CC As alternativas 'b' e 'c' estão corretas.
\end{choices}

\question[1] Como se declara uma constante em Java?
\begin{choices}
 \choice Usando a palavra-chave `const`.
 \CC Usando as palavras-chave `final` e `static`.
 \choice Apenas declarando a variável com letras maiúsculas.
 \choice Usando a palavra-chave `constant`.
 \choice Não é possível declarar constantes em Java.
\end{choices}

\question[1] O que é a "pilha de execução" (stack) em Java?
\begin{choices}
 \choice Uma estrutura de dados para armazenar objetos criados com `new`.
 \CC Uma área de memória onde são armazenadas as chamadas de métodos e variáveis locais.
 \choice Uma fila de processos esperando para serem executados pela JVM.
 \choice O local onde o bytecode é armazenado.
 \choice Uma biblioteca para operações matemáticas.
\end{choices}

\question[1] Qual a principal desvantagem de usar `==' para comparar Strings?
\begin{choices}
 \choice É mais lento que `.equals()`.
 \choice Lança uma exceção se as strings forem diferentes.
 \CC Compara as referências de memória, não o conteúdo, o que pode levar a resultados inesperados.
 \choice Não funciona para strings criadas com `new String()`.
 \choice Só funciona se as strings estiverem em letras minúsculas.
\end{choices}

\question[1] Qual é a saída do código a seguir?
\lstinputlisting{code/Ola.java}
\begin{choices}
 \choice true
 \CC false
 \choice Erro de compilação
 \choice `Olá`
 \choice Nada é impresso.
\end{choices}

\question[1] Qual é a saída do código a seguir?
\lstinputlisting{code/For.java}
\begin{choices}
 \choice 4
 \CC 5
 \choice 0
 \choice Erro de compilação, `i' fora de escopo.
 \choice 6
\end{choices}

\question[1] Qual dos seguintes tipos de dados em Java é usado para armazenar números inteiros?
\begin{choices}
 \choice double
 \choice float
 \CC int
 \choice char
 \choice boolean
\end{choices}

\question[1] Qual operador é usado para negação lógica em Java?
\begin{choices}
 \choice `&&`
 \choice `||`
 \CC `!'
 \choice `~'
 \choice `NOT'
\end{choices}

\question[1] Qual das seguintes opções é a melhor prática para definir o nome de uma classe em Java?
\begin{choices}
 \choice `minhaClasse`
 \CC `MinhaClasse`
 \choice `minha_classe`
 \choice `MINHA_CLASSE`
 \choice `minha-classe`
\end{choices}

\question[1] Qual estrutura de repetição é mais adequada quando não se sabe o número exato de iterações, mas se conhece a condição de parada?
\begin{choices}
 \choice for
 \choice foreach
 \CC while
 \choice if-else
 \choice switch
\end{choices}

% --- CONCEITOS DE POO (13 Questões) ---

\question[1] O que a anotação `@Override' indica?
\begin{choices}
 \choice Que o método está sobrecarregado.
 \CC Que o método está sendo sobrescrito de uma superclasse.
 \choice Que o método é a versão final e não pode ser modificado.
 \choice Que o método é obsoleto e não deve ser usado.
 \choice Que o método é estático.
\end{choices}

\question[1] Qual a diferença entre sobrecarga (overloading) e sobrescrita (overriding)?
\begin{choices}
 \choice Sobrecarga é em classes diferentes, sobrescrita é na mesma classe.
 \choice Sobrecarga usa a mesma assinatura, sobrescrita usa assinaturas diferentes.
 \CC Sobrecarga é ter métodos com mesmo nome e parâmetros diferentes na mesma classe; sobrescrita é redefinir um método da superclasse na subclasse.
 \choice Sobrecarga é para construtores, sobrescrita é para métodos comuns.
 \choice Não há diferença, são termos sinônimos.
\end{choices}

\question[1] Qual a diferença fundamental entre Agregação e Composição?
\begin{choices}
 \choice Na Composição, as partes podem existir sem o todo; na Agregação, não.
 \CC Na Agregação, as partes podem existir sem o todo; na Composição, as partes dependem do ciclo de vida do todo.
 \choice Agregação é uma relação ``é um'', enquanto Composição é ``tem um''.
 \choice Não há diferença, são a mesma coisa.
 \choice Composição usa a palavra-chave `composes', enquanto Agregação usa `aggregates'.
\end{choices}

\question[1] Se uma classe é declarada como `final', o que isso significa?
\begin{choices}
 \choice Que ela não pode ser instanciada.
 \choice Que todos os seus métodos são estáticos.
 \CC Que ela não pode ser herdada por outra classe.
 \choice Que ela só pode ter um único objeto.
 \choice Que seus atributos não podem ser modificados.
\end{choices}

\question[1] Qual é a saída do código a seguir?
\lstinputlisting{code/Veiculo.java}
\begin{choices}
 \choice `Veículo se move`
 \CC `Carro acelera`
 \choice Erro de compilação
 \choice `Veículo se move'`Carro acelera'
 \choice Nenhuma saída
\end{choices}

\question[1] Qual é a saída do código a seguir?
\lstinputlisting{code/Pai.java}
\begin{choices}
 \choice `Filho`
 \choice `Pai`
 \choice `Filho` `Pai`
 \CC `Pai` `Filho`
 \choice Erro de compilação
\end{choices}

\question[1] O que é encapsulamento?
\begin{choices}
 \choice A capacidade de uma classe herdar características de outra.
 \CC A prática de agrupar dados e métodos em uma unidade, escondendo os detalhes internos.
 \choice A capacidade de um objeto ter múltiplas formas.
 \choice A criação de múltiplos objetos a partir da mesma classe.
 \choice A organização de classes em pacotes.
\end{choices}

\question[1] O que é herança em POO?
\begin{choices}
 \choice A criação de uma cópia exata de um objeto.
 \CC Um mecanismo que permite a uma classe adquirir os atributos e métodos de outra.
 \choice A capacidade de um método ter o mesmo nome mas diferentes parâmetros.
 \choice A prática de esconder a complexidade interna de um objeto.
 \choice A relação onde um objeto contém outro.
\end{choices}

\question[1] Qual é a classe raiz (superclasse) de todas as classes em Java?
\begin{choices}
 \choice `Main`
 \choice `System`
 \choice `Class`
 \CC `Object`
 \choice `Java`
\end{choices}

\question[1] O que o método `toString()' da classe `Object' faz por padrão?
\begin{choices}
 \choice Retorna uma representação em String vazia do objeto.
 \CC Retorna o nome da classe seguido por "@" e o código hash do objeto.
 \choice Retorna `null`.
 \choice Lança uma exceção.
 \choice Retorna uma lista de todos os atributos e seus valores.
\end{choices}

\question[1] Uma classe `Carro' "tem um" `Motor'. Que tipo de relação de POO isso representa?
\begin{choices}
 \choice Herança
 \choice Polimorfismo
 \choice Abstração
 \CC Associação (Agregação ou Composição)
 \choice Interface
\end{choices}

\question[1] O que acontece se uma classe não define explicitamente um construtor?
\begin{choices}
 \choice A classe não pode ser instanciada.
 \choice Ocorre um erro de compilação.
 \CC O compilador Java fornece um construtor padrão (sem argumentos) automaticamente.
 \choice É necessário chamar um método `init()` para criar o objeto.
 \choice A classe só pode ter métodos estáticos.
\end{choices}

\question[1] Qual modificador de acesso torna um atributo ou método acessível apenas instâncias da própria classe?
\begin{choices}
 \choice public
 \choice protected
 \choice default
 \CC private
 \choice static
\end{choices}

% --- GREENFOOT (8 Questões) ---

\question[1] Qual método é usado para detectar se um ator está tocando em outro de uma classe específica?
\begin{choices}
 \choice `isColliding(Classe.class)`
 \choice `intersects(Classe.class)`
 \choice `hasContact(Classe.class)`
 \choice `onCollision(Classe.class)`
 \CC `isTouching(Classe.class)`
\end{choices}

\question[1] Como você removeria do mundo um objeto da classe `Moeda` que o ator atual está tocando?
\begin{choices}
 \choice `delete(Moeda.class);`
 \CC `removeTouching(Moeda.class);`
 \choice `getWorld().removeObject( getTouching( Moeda.class));`
 \choice `destroy(Moeda.class);`
 \choice `getWorld().removeTouching(Moeda.class);`
\end{choices}

\question[1] O que o método `getOneIntersectingObject(Classe.class)' retorna se não há intersecção?
\begin{choices}
 \choice `false`
 \choice `0`
 \CC `null`
 \choice Lança uma exceção `NullPointerException`.
 \choice `-1`
\end{choices}

\question[1] Qual método do Greenfoot é usado para verificar se uma tecla específica está sendo pressionada?
\begin{choices}
 \choice `Greenfoot.isKeyPressed()`
 \CC `Greenfoot.isKeyDown()`
 \choice `Greenfoot.getKey()`
 \choice `Greenfoot.checkKey()`
 \choice `Greenfoot.isPressed()`
\end{choices}

\question[1] Qual método retorna uma referência ao mundo em que o ator está?
\begin{choices}
 \choice `getParent()`
 \choice `getScene()`
 \choice `getContext()`
 \CC `getWorld()`
 \choice `getEnvironment()`
\end{choices}

\question[1] Se você quer que um ator faça algo apenas uma vez quando é adicionado ao mundo, onde você colocaria esse código?
\begin{choices}
 \choice No método `act()`, dentro de um `if`.
 \CC No construtor da classe do ator.
 \choice No construtor da classe `World`.
 \choice Na classe pai `Actor`.
\end{choices}

\question[1] Como você pode remover o próprio ator de dentro de sua classe?
\begin{choices}
 \choice `remove(this);`
 \choice `delete this;`
 \CC `getWorld().removeObject(this);`
 \choice `this.destroy();`
 \choice `Greenfoot.removeActor(this);`
\end{choices}

\question[1] Qual é o sistema de coordenadas do Greenfoot?
\begin{choices}
 \choice O ponto (0,0) está no centro do mundo.
 \choice O ponto (0,0) está no canto inferior esquerdo, com Y para cima.
 \choice O sistema de coordenadas é polar.
 \choice O sistema de coordenadas é definido pelo programador.
 \CC O ponto (0,0) está no canto superior esquerdo, com X para a direita e Y para baixo.
\end{choices}

\end{multicols}	
     \end{questions}	
\printmyanswers
\end{document}